\documentclass[aspectratio=169,10pt]{beamer}
% ---------------------------------------------------------------
% R4.02 Qualite de developpement -- preambule commun des diapositives
% Charte reprise de R3.07 (metropolis, dark tech)
% Palette IUT partagee (violet du logo + accent or)
% Utilisation :
%   \documentclass[aspectratio=169,10pt]{beamer}
%   % ---------------------------------------------------------------
% R4.02 Qualite de developpement -- preambule commun des diapositives
% Charte reprise de R3.07 (metropolis, dark tech)
% Palette IUT partagee (violet du logo + accent or)
% Utilisation :
%   \documentclass[aspectratio=169,10pt]{beamer}
%   % ---------------------------------------------------------------
% R4.02 Qualite de developpement -- preambule commun des diapositives
% Charte reprise de R3.07 (metropolis, dark tech)
% Palette IUT partagee (violet du logo + accent or)
% Utilisation :
%   \documentclass[aspectratio=169,10pt]{beamer}
%   \input{../../latex/common/r402-slides-preamble.tex}
%   \graphicspath{{../../latex/common/assets/}{img/}{figs/}}
% ---------------------------------------------------------------
\usetheme{metropolis}
\usepackage{appendixnumberbeamer}

\usepackage[T1]{fontenc}
\usepackage[utf8]{inputenc}
\usepackage[french]{babel}
\usepackage[scale=0.94]{tgheros}
\renewcommand{\familydefault}{\sfdefault}
\usepackage{amsmath,amssymb,amsfonts}
\usepackage{graphicx}
\usepackage{booktabs}
\usepackage{multirow}
\usepackage{tikz}
\usetikzlibrary{shapes.geometric, shapes, arrows, arrows.meta, positioning, calc, fit}
\usepackage{xcolor}
\usepackage{listings}
\usepackage[skins,listings]{tcolorbox}
\usepackage{qrcode}
\usepackage{eurosym}
\usepackage{hyperref}

% ---------------------------------------------------------------
% PALETTE (dark-mode tech: near-black canvas, violet (logo IUT) + gold accents)
% ---------------------------------------------------------------
\definecolor{qabg}{HTML}{0E1016}      % canvas background
\definecolor{qapanel}{HTML}{181B24}   % card / block panel
\definecolor{qapanelalt}{HTML}{20242F}% secondary panel tint
\definecolor{qaline}{HTML}{2B303D}    % hairline rules / table rules
\definecolor{qatext}{HTML}{ECEEF2}    % primary text (off-white)
\definecolor{qamuted}{HTML}{8B93A6}   % secondary / caption text
\definecolor{qaaccent}{HTML}{D93ED7}  % primary accent -- violet du logo IUT, eclairci pour fond sombre
\definecolor{qaaccent2}{HTML}{F2B84B} % secondary accent -- or/ambre
\definecolor{qacodekw}{HTML}{C792EA}  % code : mots-cles
\definecolor{qacodestr}{HTML}{C3E88D} % code : chaines
\definecolor{qacodecom}{HTML}{7F8AA3} % code : commentaires
\definecolor{qacodeann}{HTML}{F2B84B} % code : annotations / javadoc

\setbeamercolor{background canvas}{bg=qabg}
\setbeamercolor{normal text}{fg=qatext}
\setbeamercolor{palette primary}{bg=qabg, fg=qamuted}
\setbeamercolor{frametitle}{bg=qabg, fg=qatext}
\setbeamercolor{progress bar}{fg=qaaccent, bg=qaline}
\setbeamercolor{alerted text}{fg=qaaccent2}
\setbeamercolor{block title}{bg=qapanel, fg=qaaccent}
\setbeamercolor{block body}{bg=qapanel, fg=qatext}
\setbeamercolor{block title alerted}{bg=qapanel, fg=qaaccent2}
\setbeamercolor{block body alerted}{bg=qapanel, fg=qatext}
\setbeamercolor{block title example}{bg=qapanelalt, fg=qaaccent}
\setbeamercolor{block body example}{bg=qapanelalt, fg=qatext}
\setbeamercolor{example text}{fg=qaaccent}
\setbeamercolor{itemize item}{fg=qaaccent}
\setbeamercolor{itemize subitem}{fg=qaaccent2}
\setbeamercolor{enumerate item}{fg=qaaccent}
\setbeamercolor{enumerate subitem}{fg=qaaccent2}
\setbeamercolor{caption}{fg=qamuted}
\setbeamercolor{page number in head/foot}{fg=qamuted, bg=qabg}
\setbeamercolor{title}{fg=qatext}
\setbeamercolor{subtitle}{fg=qamuted}
\setbeamercolor{author}{fg=qatext}
\setbeamercolor{institute}{fg=qamuted}
\setbeamercolor{date}{fg=qamuted}
\setbeamercolor{section in toc}{fg=qaaccent}
\setbeamercolor{subsection in toc}{fg=qamuted}
% Sommaire de debut de section : section en cours en violet et gras, les autres en gris
\setbeamercolor{section in toc shaded}{fg=qamuted}
\setbeamercolor{subsection in toc shaded}{fg=qamuted}

\setbeamertemplate{blocks}[rounded][shadow=false]
\setbeamertemplate{itemize item}{\raisebox{0.15em}{\scriptsize\textcolor{qaaccent}{\ensuremath{\blacktriangleright}}}}
\setbeamertemplate{itemize subitem}{\raisebox{0.1em}{\tiny\textcolor{qaaccent2}{\ensuremath{\bullet}}}}

\setbeamerfont{frametitle}{series=\bfseries, size=\Large}
\setbeamerfont{title}{series=\bfseries, size=\Large}
\setbeamerfont{subtitle}{size=\normalsize}
\setbeamerfont{institute}{size=\scriptsize}
\setbeamerfont{date}{size=\scriptsize}
\setbeamerfont{author}{size=\small}

\metroset{sectionpage=none, subsectionpage=none, numbering=fraction, progressbar=frametitle}
\setbeamertemplate{section in toc}[sections numbered]

% ---------------------------------------------------------------
% LISTINGS -- theme sombre (lstlisting reste utilisable tel quel)
% \lstcode : police + couleur de base, a reutiliser dans les
%            surcharges locales (ex. basicstyle=\lstcode\scriptsize)
% ---------------------------------------------------------------
\newcommand{\lstcode}{\ttfamily\color{qatext}}

\lstdefinelanguage{TypeScript}{
  keywords={typeof, new, true, false, catch, function, return, null, switch, var,
    if, in, while, do, else, case, break, const, let, interface, class, extends,
    implements, enum, export, import, this, super, constructor, public, private,
    protected, static, async, await, throw, throws, try, finally, void},
  ndkeywords={boolean, any, number, string, int, long, double},
  sensitive=true,
  comment=[l]{//},
  morecomment=[s]{/*}{*/},
  morestring=[b]',
  morestring=[b]"
}

\lstset{
  language=Java,
  basicstyle=\lstcode\footnotesize,
  keywordstyle=\color{qacodekw}\bfseries,
  ndkeywordstyle=\color{qaaccent2}\bfseries,
  identifierstyle=\color{qatext},
  stringstyle=\color{qacodestr},
  commentstyle=\color{qacodecom}\itshape,
  morecomment=[s][\color{qacodeann}]{/**}{*/},
  backgroundcolor=\color{qapanelalt},
  rulecolor=\color{qaline},
  frame=single,
  framesep=3pt,
  xleftmargin=4pt,
  xrightmargin=4pt,
  keepspaces=true,
  columns=fullflexible,
  showstringspaces=false,
  breaklines=true,
  tabsize=2,
  escapeinside={(*@}{@*)},
  % accents UTF-8 dans le code (listings ne gere pas l'UTF-8 multi-octets)
  literate={é}{{\'e}}1 {è}{{\`e}}1 {ê}{{\^e}}1 {à}{{\`a}}1 {â}{{\^a}}1
           {ù}{{\`u}}1 {û}{{\^u}}1 {ô}{{\^o}}1 {î}{{\^i}}1 {ç}{{\c c}}1
           {É}{{\'E}}1 {À}{{\`A}}1 {€}{{\euro{}}}1 {…}{{\ldots}}1,
}

% ---------------------------------------------------------------
% CARD-STYLE BOXES
% ---------------------------------------------------------------
\tcbset{
  figcard/.style={colback=white, colframe=white, boxrule=0pt, arc=4pt,
    left=3pt, right=3pt, top=3pt, bottom=3pt, boxsep=0pt,
    halign=center, valign=center},
  datacard/.style={colback=qapanel, colframe=qaline, boxrule=0.4pt, arc=3pt,
    left=8pt, right=8pt, top=6pt, bottom=6pt, coltext=qatext,
    before upper={\color{qatext}}, halign=center},
  codebox/.style={colback=qapanelalt, colframe=qaline, boxrule=0.4pt, arc=3pt,
    left=8pt, right=8pt, top=5pt, bottom=5pt, listing only,
    listing options={basicstyle=\ttfamily\small\color{qaaccent},
      breaklines=true, showstringspaces=false, columns=fullflexible,
      keepspaces=true, frame=none, backgroundcolor={}}}
}
\newtcolorbox{figcard}[1][]{figcard,#1}
\newtcolorbox{datacard}{datacard}
\newtcblisting{codebox}{codebox}

% ---------------------------------------------------------------
% HANDY MACROS
% ---------------------------------------------------------------
% Mention d'elaboration (discrete : plus sombre que le texte secondaire)
\definecolor{qaaimention}{HTML}{4E5566}
\newcommand{\raimention}{{\tiny\color{qaaimention}Supports \'elabor\'es avec l'aide de Claude Code Opus 5.5}}
\newcommand{\hl}[1]{\alert{#1}}
\newcommand{\source}[1]{{\scriptsize\color{qamuted}#1}}

% Sommaire automatique en debut de section
\AtBeginSection[]{%
  \begin{frame}<beamer>{Plan du cours}
    \tableofcontents[currentsection]
  \end{frame}}

% Page de titre : \rtitleslide{numero du cours}{titre du cours}
% Les auteurs sont repris de \rauthors (modifiable avant \begin{document})
\newcommand{\rauthors}{Chouki Tibermacine}
\newcommand{\rtitleslide}[2]{%
\begin{frame}[plain,noframenumbering]
  \begin{center}
    \vspace{2mm}
    {\color{qaaccent}\rule{2.1cm}{1.2pt}}\\[6mm]
    {\Large\bfseries\color{qatext} R4.02}\\[1mm]
    {\Large\bfseries\color{qatext} Qualit\'e de d\'eveloppement}\\[5mm]
    {\large\color{qamuted} Cours #1}\\[1mm]
    {\large\color{qaaccent} #2}\\[6mm]
    \includegraphics[height=2.1cm]{iut_logo_transparent.png}

    \vspace{4mm}

    {\small\color{qatext} \rauthors}\\[1mm]
    {\scriptsize\color{qamuted} BUT 2 Informatique --- IUT de Vannes --- Universit\'e Bretagne Sud}\\[2mm]
    \raimention
  \end{center}
\end{frame}
}

% ---- Diapo de demonstration : \rdemo{numero}{titre}{contenu (items)}{fichier(s) a ouvrir} ----
\newcommand{\rdemo}[4]{%
\begin{frame}{D\'emo #1 --- #2}
  \begin{tikzpicture}[remember picture, overlay]
    \node[anchor=north east, fill=qaaccent2, text=qabg, font=\scriptsize\bfseries,
      rounded corners=3pt, inner sep=4pt] at ([xshift=-0.4cm, yshift=-0.35cm]current page.north east) {D\'EMO};
  \end{tikzpicture}
  \begin{columns}[T]
    \begin{column}{0.62\textwidth}
      \begin{itemize}
        #3
      \end{itemize}
    \end{column}
    \begin{column}{0.34\textwidth}
      \begin{datacard}
        \raggedright\scriptsize
        {\color{qamuted} Dans l'IDE}\\[1mm]
        #4
      \end{datacard}
    \end{column}
  \end{columns}
\end{frame}}

% ---- Diapo de fin « Questions ? » ----
\newcommand{\rcontact}{chouki.tibermacine@univ-ubs.fr}
\newcommand{\rquestionslide}{%
\begin{frame}{Questions ?}
  \begin{center}
    \vspace{2mm}
    \begin{tikzpicture}
      \node[regular polygon, regular polygon sides=6, draw=qaline, thick,
        minimum size=4.6cm, rotate=90] (hexin) {};
      \node[regular polygon, regular polygon sides=6, draw=qaaccent, dashed,
        thick, minimum size=6.2cm, rotate=90] (hexout) {};
      \node[font=\fontsize{92}{92}\selectfont\bfseries, color=qaaccent2] at (hexin) {?};
    \end{tikzpicture}\\[5mm]
    {\small\color{qamuted} \rcontact}
  \end{center}
\end{frame}}

%   \graphicspath{{../../latex/common/assets/}{img/}{figs/}}
% ---------------------------------------------------------------
\usetheme{metropolis}
\usepackage{appendixnumberbeamer}

\usepackage[T1]{fontenc}
\usepackage[utf8]{inputenc}
\usepackage[french]{babel}
\usepackage[scale=0.94]{tgheros}
\renewcommand{\familydefault}{\sfdefault}
\usepackage{amsmath,amssymb,amsfonts}
\usepackage{graphicx}
\usepackage{booktabs}
\usepackage{multirow}
\usepackage{tikz}
\usetikzlibrary{shapes.geometric, shapes, arrows, arrows.meta, positioning, calc, fit}
\usepackage{xcolor}
\usepackage{listings}
\usepackage[skins,listings]{tcolorbox}
\usepackage{qrcode}
\usepackage{eurosym}
\usepackage{hyperref}

% ---------------------------------------------------------------
% PALETTE (dark-mode tech: near-black canvas, violet (logo IUT) + gold accents)
% ---------------------------------------------------------------
\definecolor{qabg}{HTML}{0E1016}      % canvas background
\definecolor{qapanel}{HTML}{181B24}   % card / block panel
\definecolor{qapanelalt}{HTML}{20242F}% secondary panel tint
\definecolor{qaline}{HTML}{2B303D}    % hairline rules / table rules
\definecolor{qatext}{HTML}{ECEEF2}    % primary text (off-white)
\definecolor{qamuted}{HTML}{8B93A6}   % secondary / caption text
\definecolor{qaaccent}{HTML}{D93ED7}  % primary accent -- violet du logo IUT, eclairci pour fond sombre
\definecolor{qaaccent2}{HTML}{F2B84B} % secondary accent -- or/ambre
\definecolor{qacodekw}{HTML}{C792EA}  % code : mots-cles
\definecolor{qacodestr}{HTML}{C3E88D} % code : chaines
\definecolor{qacodecom}{HTML}{7F8AA3} % code : commentaires
\definecolor{qacodeann}{HTML}{F2B84B} % code : annotations / javadoc

\setbeamercolor{background canvas}{bg=qabg}
\setbeamercolor{normal text}{fg=qatext}
\setbeamercolor{palette primary}{bg=qabg, fg=qamuted}
\setbeamercolor{frametitle}{bg=qabg, fg=qatext}
\setbeamercolor{progress bar}{fg=qaaccent, bg=qaline}
\setbeamercolor{alerted text}{fg=qaaccent2}
\setbeamercolor{block title}{bg=qapanel, fg=qaaccent}
\setbeamercolor{block body}{bg=qapanel, fg=qatext}
\setbeamercolor{block title alerted}{bg=qapanel, fg=qaaccent2}
\setbeamercolor{block body alerted}{bg=qapanel, fg=qatext}
\setbeamercolor{block title example}{bg=qapanelalt, fg=qaaccent}
\setbeamercolor{block body example}{bg=qapanelalt, fg=qatext}
\setbeamercolor{example text}{fg=qaaccent}
\setbeamercolor{itemize item}{fg=qaaccent}
\setbeamercolor{itemize subitem}{fg=qaaccent2}
\setbeamercolor{enumerate item}{fg=qaaccent}
\setbeamercolor{enumerate subitem}{fg=qaaccent2}
\setbeamercolor{caption}{fg=qamuted}
\setbeamercolor{page number in head/foot}{fg=qamuted, bg=qabg}
\setbeamercolor{title}{fg=qatext}
\setbeamercolor{subtitle}{fg=qamuted}
\setbeamercolor{author}{fg=qatext}
\setbeamercolor{institute}{fg=qamuted}
\setbeamercolor{date}{fg=qamuted}
\setbeamercolor{section in toc}{fg=qaaccent}
\setbeamercolor{subsection in toc}{fg=qamuted}
% Sommaire de debut de section : section en cours en violet et gras, les autres en gris
\setbeamercolor{section in toc shaded}{fg=qamuted}
\setbeamercolor{subsection in toc shaded}{fg=qamuted}

\setbeamertemplate{blocks}[rounded][shadow=false]
\setbeamertemplate{itemize item}{\raisebox{0.15em}{\scriptsize\textcolor{qaaccent}{\ensuremath{\blacktriangleright}}}}
\setbeamertemplate{itemize subitem}{\raisebox{0.1em}{\tiny\textcolor{qaaccent2}{\ensuremath{\bullet}}}}

\setbeamerfont{frametitle}{series=\bfseries, size=\Large}
\setbeamerfont{title}{series=\bfseries, size=\Large}
\setbeamerfont{subtitle}{size=\normalsize}
\setbeamerfont{institute}{size=\scriptsize}
\setbeamerfont{date}{size=\scriptsize}
\setbeamerfont{author}{size=\small}

\metroset{sectionpage=none, subsectionpage=none, numbering=fraction, progressbar=frametitle}
\setbeamertemplate{section in toc}[sections numbered]

% ---------------------------------------------------------------
% LISTINGS -- theme sombre (lstlisting reste utilisable tel quel)
% \lstcode : police + couleur de base, a reutiliser dans les
%            surcharges locales (ex. basicstyle=\lstcode\scriptsize)
% ---------------------------------------------------------------
\newcommand{\lstcode}{\ttfamily\color{qatext}}

\lstdefinelanguage{TypeScript}{
  keywords={typeof, new, true, false, catch, function, return, null, switch, var,
    if, in, while, do, else, case, break, const, let, interface, class, extends,
    implements, enum, export, import, this, super, constructor, public, private,
    protected, static, async, await, throw, throws, try, finally, void},
  ndkeywords={boolean, any, number, string, int, long, double},
  sensitive=true,
  comment=[l]{//},
  morecomment=[s]{/*}{*/},
  morestring=[b]',
  morestring=[b]"
}

\lstset{
  language=Java,
  basicstyle=\lstcode\footnotesize,
  keywordstyle=\color{qacodekw}\bfseries,
  ndkeywordstyle=\color{qaaccent2}\bfseries,
  identifierstyle=\color{qatext},
  stringstyle=\color{qacodestr},
  commentstyle=\color{qacodecom}\itshape,
  morecomment=[s][\color{qacodeann}]{/**}{*/},
  backgroundcolor=\color{qapanelalt},
  rulecolor=\color{qaline},
  frame=single,
  framesep=3pt,
  xleftmargin=4pt,
  xrightmargin=4pt,
  keepspaces=true,
  columns=fullflexible,
  showstringspaces=false,
  breaklines=true,
  tabsize=2,
  escapeinside={(*@}{@*)},
  % accents UTF-8 dans le code (listings ne gere pas l'UTF-8 multi-octets)
  literate={é}{{\'e}}1 {è}{{\`e}}1 {ê}{{\^e}}1 {à}{{\`a}}1 {â}{{\^a}}1
           {ù}{{\`u}}1 {û}{{\^u}}1 {ô}{{\^o}}1 {î}{{\^i}}1 {ç}{{\c c}}1
           {É}{{\'E}}1 {À}{{\`A}}1 {€}{{\euro{}}}1 {…}{{\ldots}}1,
}

% ---------------------------------------------------------------
% CARD-STYLE BOXES
% ---------------------------------------------------------------
\tcbset{
  figcard/.style={colback=white, colframe=white, boxrule=0pt, arc=4pt,
    left=3pt, right=3pt, top=3pt, bottom=3pt, boxsep=0pt,
    halign=center, valign=center},
  datacard/.style={colback=qapanel, colframe=qaline, boxrule=0.4pt, arc=3pt,
    left=8pt, right=8pt, top=6pt, bottom=6pt, coltext=qatext,
    before upper={\color{qatext}}, halign=center},
  codebox/.style={colback=qapanelalt, colframe=qaline, boxrule=0.4pt, arc=3pt,
    left=8pt, right=8pt, top=5pt, bottom=5pt, listing only,
    listing options={basicstyle=\ttfamily\small\color{qaaccent},
      breaklines=true, showstringspaces=false, columns=fullflexible,
      keepspaces=true, frame=none, backgroundcolor={}}}
}
\newtcolorbox{figcard}[1][]{figcard,#1}
\newtcolorbox{datacard}{datacard}
\newtcblisting{codebox}{codebox}

% ---------------------------------------------------------------
% HANDY MACROS
% ---------------------------------------------------------------
% Mention d'elaboration (discrete : plus sombre que le texte secondaire)
\definecolor{qaaimention}{HTML}{4E5566}
\newcommand{\raimention}{{\tiny\color{qaaimention}Supports \'elabor\'es avec l'aide de Claude Code Opus 5.5}}
\newcommand{\hl}[1]{\alert{#1}}
\newcommand{\source}[1]{{\scriptsize\color{qamuted}#1}}

% Sommaire automatique en debut de section
\AtBeginSection[]{%
  \begin{frame}<beamer>{Plan du cours}
    \tableofcontents[currentsection]
  \end{frame}}

% Page de titre : \rtitleslide{numero du cours}{titre du cours}
% Les auteurs sont repris de \rauthors (modifiable avant \begin{document})
\newcommand{\rauthors}{Chouki Tibermacine}
\newcommand{\rtitleslide}[2]{%
\begin{frame}[plain,noframenumbering]
  \begin{center}
    \vspace{2mm}
    {\color{qaaccent}\rule{2.1cm}{1.2pt}}\\[6mm]
    {\Large\bfseries\color{qatext} R4.02}\\[1mm]
    {\Large\bfseries\color{qatext} Qualit\'e de d\'eveloppement}\\[5mm]
    {\large\color{qamuted} Cours #1}\\[1mm]
    {\large\color{qaaccent} #2}\\[6mm]
    \includegraphics[height=2.1cm]{iut_logo_transparent.png}

    \vspace{4mm}

    {\small\color{qatext} \rauthors}\\[1mm]
    {\scriptsize\color{qamuted} BUT 2 Informatique --- IUT de Vannes --- Universit\'e Bretagne Sud}\\[2mm]
    \raimention
  \end{center}
\end{frame}
}

% ---- Diapo de demonstration : \rdemo{numero}{titre}{contenu (items)}{fichier(s) a ouvrir} ----
\newcommand{\rdemo}[4]{%
\begin{frame}{D\'emo #1 --- #2}
  \begin{tikzpicture}[remember picture, overlay]
    \node[anchor=north east, fill=qaaccent2, text=qabg, font=\scriptsize\bfseries,
      rounded corners=3pt, inner sep=4pt] at ([xshift=-0.4cm, yshift=-0.35cm]current page.north east) {D\'EMO};
  \end{tikzpicture}
  \begin{columns}[T]
    \begin{column}{0.62\textwidth}
      \begin{itemize}
        #3
      \end{itemize}
    \end{column}
    \begin{column}{0.34\textwidth}
      \begin{datacard}
        \raggedright\scriptsize
        {\color{qamuted} Dans l'IDE}\\[1mm]
        #4
      \end{datacard}
    \end{column}
  \end{columns}
\end{frame}}

% ---- Diapo de fin « Questions ? » ----
\newcommand{\rcontact}{chouki.tibermacine@univ-ubs.fr}
\newcommand{\rquestionslide}{%
\begin{frame}{Questions ?}
  \begin{center}
    \vspace{2mm}
    \begin{tikzpicture}
      \node[regular polygon, regular polygon sides=6, draw=qaline, thick,
        minimum size=4.6cm, rotate=90] (hexin) {};
      \node[regular polygon, regular polygon sides=6, draw=qaaccent, dashed,
        thick, minimum size=6.2cm, rotate=90] (hexout) {};
      \node[font=\fontsize{92}{92}\selectfont\bfseries, color=qaaccent2] at (hexin) {?};
    \end{tikzpicture}\\[5mm]
    {\small\color{qamuted} \rcontact}
  \end{center}
\end{frame}}

%   \graphicspath{{../../latex/common/assets/}{img/}{figs/}}
% ---------------------------------------------------------------
\usetheme{metropolis}
\usepackage{appendixnumberbeamer}

\usepackage[T1]{fontenc}
\usepackage[utf8]{inputenc}
\usepackage[french]{babel}
\usepackage[scale=0.94]{tgheros}
\renewcommand{\familydefault}{\sfdefault}
\usepackage{amsmath,amssymb,amsfonts}
\usepackage{graphicx}
\usepackage{booktabs}
\usepackage{multirow}
\usepackage{tikz}
\usetikzlibrary{shapes.geometric, shapes, arrows, arrows.meta, positioning, calc, fit}
\usepackage{xcolor}
\usepackage{listings}
\usepackage[skins,listings]{tcolorbox}
\usepackage{qrcode}
\usepackage{eurosym}
\usepackage{hyperref}

% ---------------------------------------------------------------
% PALETTE (dark-mode tech: near-black canvas, violet (logo IUT) + gold accents)
% ---------------------------------------------------------------
\definecolor{qabg}{HTML}{0E1016}      % canvas background
\definecolor{qapanel}{HTML}{181B24}   % card / block panel
\definecolor{qapanelalt}{HTML}{20242F}% secondary panel tint
\definecolor{qaline}{HTML}{2B303D}    % hairline rules / table rules
\definecolor{qatext}{HTML}{ECEEF2}    % primary text (off-white)
\definecolor{qamuted}{HTML}{8B93A6}   % secondary / caption text
\definecolor{qaaccent}{HTML}{D93ED7}  % primary accent -- violet du logo IUT, eclairci pour fond sombre
\definecolor{qaaccent2}{HTML}{F2B84B} % secondary accent -- or/ambre
\definecolor{qacodekw}{HTML}{C792EA}  % code : mots-cles
\definecolor{qacodestr}{HTML}{C3E88D} % code : chaines
\definecolor{qacodecom}{HTML}{7F8AA3} % code : commentaires
\definecolor{qacodeann}{HTML}{F2B84B} % code : annotations / javadoc

\setbeamercolor{background canvas}{bg=qabg}
\setbeamercolor{normal text}{fg=qatext}
\setbeamercolor{palette primary}{bg=qabg, fg=qamuted}
\setbeamercolor{frametitle}{bg=qabg, fg=qatext}
\setbeamercolor{progress bar}{fg=qaaccent, bg=qaline}
\setbeamercolor{alerted text}{fg=qaaccent2}
\setbeamercolor{block title}{bg=qapanel, fg=qaaccent}
\setbeamercolor{block body}{bg=qapanel, fg=qatext}
\setbeamercolor{block title alerted}{bg=qapanel, fg=qaaccent2}
\setbeamercolor{block body alerted}{bg=qapanel, fg=qatext}
\setbeamercolor{block title example}{bg=qapanelalt, fg=qaaccent}
\setbeamercolor{block body example}{bg=qapanelalt, fg=qatext}
\setbeamercolor{example text}{fg=qaaccent}
\setbeamercolor{itemize item}{fg=qaaccent}
\setbeamercolor{itemize subitem}{fg=qaaccent2}
\setbeamercolor{enumerate item}{fg=qaaccent}
\setbeamercolor{enumerate subitem}{fg=qaaccent2}
\setbeamercolor{caption}{fg=qamuted}
\setbeamercolor{page number in head/foot}{fg=qamuted, bg=qabg}
\setbeamercolor{title}{fg=qatext}
\setbeamercolor{subtitle}{fg=qamuted}
\setbeamercolor{author}{fg=qatext}
\setbeamercolor{institute}{fg=qamuted}
\setbeamercolor{date}{fg=qamuted}
\setbeamercolor{section in toc}{fg=qaaccent}
\setbeamercolor{subsection in toc}{fg=qamuted}
% Sommaire de debut de section : section en cours en violet et gras, les autres en gris
\setbeamercolor{section in toc shaded}{fg=qamuted}
\setbeamercolor{subsection in toc shaded}{fg=qamuted}

\setbeamertemplate{blocks}[rounded][shadow=false]
\setbeamertemplate{itemize item}{\raisebox{0.15em}{\scriptsize\textcolor{qaaccent}{\ensuremath{\blacktriangleright}}}}
\setbeamertemplate{itemize subitem}{\raisebox{0.1em}{\tiny\textcolor{qaaccent2}{\ensuremath{\bullet}}}}

\setbeamerfont{frametitle}{series=\bfseries, size=\Large}
\setbeamerfont{title}{series=\bfseries, size=\Large}
\setbeamerfont{subtitle}{size=\normalsize}
\setbeamerfont{institute}{size=\scriptsize}
\setbeamerfont{date}{size=\scriptsize}
\setbeamerfont{author}{size=\small}

\metroset{sectionpage=none, subsectionpage=none, numbering=fraction, progressbar=frametitle}
\setbeamertemplate{section in toc}[sections numbered]

% ---------------------------------------------------------------
% LISTINGS -- theme sombre (lstlisting reste utilisable tel quel)
% \lstcode : police + couleur de base, a reutiliser dans les
%            surcharges locales (ex. basicstyle=\lstcode\scriptsize)
% ---------------------------------------------------------------
\newcommand{\lstcode}{\ttfamily\color{qatext}}

\lstdefinelanguage{TypeScript}{
  keywords={typeof, new, true, false, catch, function, return, null, switch, var,
    if, in, while, do, else, case, break, const, let, interface, class, extends,
    implements, enum, export, import, this, super, constructor, public, private,
    protected, static, async, await, throw, throws, try, finally, void},
  ndkeywords={boolean, any, number, string, int, long, double},
  sensitive=true,
  comment=[l]{//},
  morecomment=[s]{/*}{*/},
  morestring=[b]',
  morestring=[b]"
}

\lstset{
  language=Java,
  basicstyle=\lstcode\footnotesize,
  keywordstyle=\color{qacodekw}\bfseries,
  ndkeywordstyle=\color{qaaccent2}\bfseries,
  identifierstyle=\color{qatext},
  stringstyle=\color{qacodestr},
  commentstyle=\color{qacodecom}\itshape,
  morecomment=[s][\color{qacodeann}]{/**}{*/},
  backgroundcolor=\color{qapanelalt},
  rulecolor=\color{qaline},
  frame=single,
  framesep=3pt,
  xleftmargin=4pt,
  xrightmargin=4pt,
  keepspaces=true,
  columns=fullflexible,
  showstringspaces=false,
  breaklines=true,
  tabsize=2,
  escapeinside={(*@}{@*)},
  % accents UTF-8 dans le code (listings ne gere pas l'UTF-8 multi-octets)
  literate={é}{{\'e}}1 {è}{{\`e}}1 {ê}{{\^e}}1 {à}{{\`a}}1 {â}{{\^a}}1
           {ù}{{\`u}}1 {û}{{\^u}}1 {ô}{{\^o}}1 {î}{{\^i}}1 {ç}{{\c c}}1
           {É}{{\'E}}1 {À}{{\`A}}1 {€}{{\euro{}}}1 {…}{{\ldots}}1,
}

% ---------------------------------------------------------------
% CARD-STYLE BOXES
% ---------------------------------------------------------------
\tcbset{
  figcard/.style={colback=white, colframe=white, boxrule=0pt, arc=4pt,
    left=3pt, right=3pt, top=3pt, bottom=3pt, boxsep=0pt,
    halign=center, valign=center},
  datacard/.style={colback=qapanel, colframe=qaline, boxrule=0.4pt, arc=3pt,
    left=8pt, right=8pt, top=6pt, bottom=6pt, coltext=qatext,
    before upper={\color{qatext}}, halign=center},
  codebox/.style={colback=qapanelalt, colframe=qaline, boxrule=0.4pt, arc=3pt,
    left=8pt, right=8pt, top=5pt, bottom=5pt, listing only,
    listing options={basicstyle=\ttfamily\small\color{qaaccent},
      breaklines=true, showstringspaces=false, columns=fullflexible,
      keepspaces=true, frame=none, backgroundcolor={}}}
}
\newtcolorbox{figcard}[1][]{figcard,#1}
\newtcolorbox{datacard}{datacard}
\newtcblisting{codebox}{codebox}

% ---------------------------------------------------------------
% HANDY MACROS
% ---------------------------------------------------------------
% Mention d'elaboration (discrete : plus sombre que le texte secondaire)
\definecolor{qaaimention}{HTML}{4E5566}
\newcommand{\raimention}{{\tiny\color{qaaimention}Supports \'elabor\'es avec l'aide de Claude Code Opus 5.5}}
\newcommand{\hl}[1]{\alert{#1}}
\newcommand{\source}[1]{{\scriptsize\color{qamuted}#1}}

% Sommaire automatique en debut de section
\AtBeginSection[]{%
  \begin{frame}<beamer>{Plan du cours}
    \tableofcontents[currentsection]
  \end{frame}}

% Page de titre : \rtitleslide{numero du cours}{titre du cours}
% Les auteurs sont repris de \rauthors (modifiable avant \begin{document})
\newcommand{\rauthors}{Chouki Tibermacine}
\newcommand{\rtitleslide}[2]{%
\begin{frame}[plain,noframenumbering]
  \begin{center}
    \vspace{2mm}
    {\color{qaaccent}\rule{2.1cm}{1.2pt}}\\[6mm]
    {\Large\bfseries\color{qatext} R4.02}\\[1mm]
    {\Large\bfseries\color{qatext} Qualit\'e de d\'eveloppement}\\[5mm]
    {\large\color{qamuted} Cours #1}\\[1mm]
    {\large\color{qaaccent} #2}\\[6mm]
    \includegraphics[height=2.1cm]{iut_logo_transparent.png}

    \vspace{4mm}

    {\small\color{qatext} \rauthors}\\[1mm]
    {\scriptsize\color{qamuted} BUT 2 Informatique --- IUT de Vannes --- Universit\'e Bretagne Sud}\\[2mm]
    \raimention
  \end{center}
\end{frame}
}

% ---- Diapo de demonstration : \rdemo{numero}{titre}{contenu (items)}{fichier(s) a ouvrir} ----
\newcommand{\rdemo}[4]{%
\begin{frame}{D\'emo #1 --- #2}
  \begin{tikzpicture}[remember picture, overlay]
    \node[anchor=north east, fill=qaaccent2, text=qabg, font=\scriptsize\bfseries,
      rounded corners=3pt, inner sep=4pt] at ([xshift=-0.4cm, yshift=-0.35cm]current page.north east) {D\'EMO};
  \end{tikzpicture}
  \begin{columns}[T]
    \begin{column}{0.62\textwidth}
      \begin{itemize}
        #3
      \end{itemize}
    \end{column}
    \begin{column}{0.34\textwidth}
      \begin{datacard}
        \raggedright\scriptsize
        {\color{qamuted} Dans l'IDE}\\[1mm]
        #4
      \end{datacard}
    \end{column}
  \end{columns}
\end{frame}}

% ---- Diapo de fin « Questions ? » ----
\newcommand{\rcontact}{chouki.tibermacine@univ-ubs.fr}
\newcommand{\rquestionslide}{%
\begin{frame}{Questions ?}
  \begin{center}
    \vspace{2mm}
    \begin{tikzpicture}
      \node[regular polygon, regular polygon sides=6, draw=qaline, thick,
        minimum size=4.6cm, rotate=90] (hexin) {};
      \node[regular polygon, regular polygon sides=6, draw=qaaccent, dashed,
        thick, minimum size=6.2cm, rotate=90] (hexout) {};
      \node[font=\fontsize{92}{92}\selectfont\bfseries, color=qaaccent2] at (hexin) {?};
    \end{tikzpicture}\\[5mm]
    {\small\color{qamuted} \rcontact}
  \end{center}
\end{frame}}

\graphicspath{{../../latex/common/assets/}{img/}}

\title{Tests de bout en bout et tests d'interface}
\author[C. Tibermacine]{Chouki Tibermacine}
\date{}

\tikzset{
  comp/.style={rectangle, rounded corners=2pt, draw=qaaccent, thick, fill=qapanel, text=qatext,
    minimum height=0.6cm, minimum width=2.3cm, align=center, font=\scriptsize},
  arr/.style={-Latex, qamuted, thick}
}

\begin{document}

\rtitleslide{5}{Tests de bout en bout et tests d'interface}

\begin{frame}{Où en est-on ?}
  \begin{columns}[c]
    \begin{column}{0.5\textwidth}
      \begin{itemize}
        \item CM1, CM3 : la logique métier, en unitaire
        \item CM2 : les critères d'acceptation, exécutables
        \item CM4 : le serveur entier, jusqu'à la base
        \item Il reste une partie de l'application qu'aucun test n'a vue : \textbf{le navigateur}
          \begin{itemize}
            \item le code JavaScript du front Vue
            \item le dialogue entre le front et l'API
            \item ce que l'utilisateur voit et fait réellement
          \end{itemize}
      \end{itemize}
    \end{column}
    \begin{column}{0.46\textwidth}
      \begin{center}
      \begin{tikzpicture}[x=0.8cm, y=1cm]
        \fill[qaaccent!35!qapanel] (0,0) -- (6,0) -- (4.6,1.4) -- (1.4,1.4) -- cycle;
        \fill[qaaccent!35!qapanel] (1.4,1.4) -- (4.6,1.4) -- (3.7,2.3) -- (2.3,2.3) -- cycle;
        \fill[qaaccent2] (2.3,2.3) -- (3.7,2.3) -- (3,3.0) -- cycle;
        \node[font=\small\bfseries, text=qatext] at (3,0.7) {Unitaires};
        \node[font=\scriptsize\bfseries, text=qatext] at (3,1.85) {Intégration};
        \node[font=\tiny\bfseries, text=qabg] at (3,2.55) {E2E};
        \draw[arr, qaaccent2] (5.4,2.6) -- (3.8,2.6);
        \node[font=\scriptsize, text=qaaccent2, anchor=west] at (5.4,2.6) {CM5};
      \end{tikzpicture}
      \end{center}
    \end{column}
  \end{columns}
\end{frame}

% ================================================================
\section{Les tests de bout en bout}

\begin{frame}{Un test de bout en bout}
  \begin{center}
  \begin{tikzpicture}[node distance=0.55cm]
    \node[comp, draw=qaaccent2] (test) {Test\\{\tiny\color{qamuted}Playwright}};
    \node[comp, right=of test] (nav) {Navigateur\\{\tiny\color{qamuted}Chromium, Firefox\ldots}};
    \node[comp, right=of nav] (front) {Front Vue\\{\tiny\color{qamuted}JavaScript}};
    \node[comp, right=of front] (api) {API REST\\{\tiny\color{qamuted}Spring Boot}};
    \node[comp, right=of api] (db) {PostgreSQL};
    \draw[arr, qaaccent2] (test) -- node[above, font=\tiny, text=qamuted] {pilote} (nav);
    \draw[arr] (nav) -- (front);
    \draw[arr] (front) -- node[above, font=\tiny, text=qamuted] {HTTP} (api);
    \draw[arr] (api) -- (db);
  \end{tikzpicture}
  \end{center}
  \vspace{2mm}
  \begin{itemize}
    \item Le test \textbf{pilote un vrai navigateur}, comme un utilisateur : il ouvre la page, remplit un formulaire, clique, lit ce qui s'affiche
    \item Toute l'application est réelle : aucune doublure (sauf, parfois, les services externes payants)
    \item Il vérifie un \textbf{parcours} : « se connecter, faire un virement, voir le nouveau solde »
    \item Il est souvent fonctionnel et boîte noire : il ne connaît que ce qui est à l'écran
  \end{itemize}
\end{frame}

\begin{frame}{Le prix à payer}
  \begin{columns}[T]
    \begin{column}{0.48\textwidth}
      \begin{alertblock}{Coûteux}
        \begin{itemize}
          \item \textbf{Lents} : démarrer l'application et un navigateur, attendre chaque page (secondes par test)
          \item \textbf{Difficiles à diagnostiquer} : « le solde affiché est faux », mais à cause de quelle couche ?
          \item \textbf{Fragiles} : un libellé ou une mise en page change, le test casse
        \end{itemize}
      \end{alertblock}
    \end{column}
    \begin{column}{0.48\textwidth}
      \begin{exampleblock}{Donc peu nombreux et bien choisis}
        \begin{itemize}
          \item Les \textbf{parcours critiques} : ceux qui coûtent cher s'ils cassent
          \item Pour la banque : se connecter, consulter ses comptes, faire un virement, être averti d'un refus
          \item Pas un test par règle métier : les règles sont déjà testées plus bas dans la pyramide
        \end{itemize}
      \end{exampleblock}
    \end{column}
  \end{columns}
\end{frame}

\begin{frame}{Les tests instables (\textit{flaky tests})}
  \begin{itemize}
    \item Un test \textbf{instable} passe ou échoue \textbf{sans que le code ait changé}
    \item Causes fréquentes en bout en bout :
      \begin{itemize}
        \item le \textbf{temps} : on lit la page avant que le serveur ait répondu, avant la fin d'une animation
        \item les \textbf{données} : un test précédent a modifié les comptes
        \item l'\textbf{environnement} : machine de CI plus lente, réseau, fuseau horaire, langue du navigateur
      \end{itemize}
    \item Conséquence : l'équipe prend l'habitude de « relancer jusqu'à ce que ça passe »\ldots
      \begin{itemize}
        \item et ne croit plus aux tests, même quand ils signalent un vrai bug
      \end{itemize}
  \end{itemize}
  \vspace{2mm}
  \begin{datacard}
    \small Un test instable n'est pas un test « un peu rouge » : c'est un test \alert{à réparer en priorité}, ou à supprimer.
  \end{datacard}
\end{frame}

\begin{frame}{Trois outils libres et gratuits}
  \footnotesize
  \begin{datacard}
    \begin{tabular}{@{}p{2.3cm}p{3.1cm}p{3.1cm}p{3.3cm}@{}}
      \toprule
      & \textbf{Selenium} & \textbf{Cypress} & \textbf{Playwright} \\
      \midrule
      Depuis & 2004 & 2017 & 2020 (Microsoft) \\[0.5mm]
      Langages & Java, Python, C\#, JS\ldots & JavaScript / TypeScript & Java, Python, C\#, JS / TS \\[0.5mm]
      Navigateurs & tous (WebDriver) & Chromium, Firefox & Chromium, Firefox, WebKit \\[0.5mm]
      Attente automatique & non : à écrire soi-même & oui & oui \\[0.5mm]
      Outils & Selenium IDE, Grid & exécution interactive & générateur de code, \textit{trace viewer} \\[0.5mm]
      Aujourd'hui & très présent dans l'existant & stable, en perte de vitesse & \textbf{le plus adopté} pour les nouveaux projets \\
      \bottomrule
    \end{tabular}
  \end{datacard}
  \vspace{2mm}
  \small
  \begin{itemize}
    \item Ici : \textbf{Playwright for Java}, dans le même projet Maven, avec JUnit, AssertJ et Cucumber
    \item Les principes de ce cours (sélecteurs, attentes, Page Object) valent pour les trois outils
  \end{itemize}
\end{frame}

% ================================================================
\section{Écrire des tests d'interface robustes}

\begin{frame}[fragile]{Un premier test Playwright}
  \begin{columns}[T]
    \begin{column}{0.64\textwidth}
\begin{lstlisting}[basicstyle=\lstcode\scriptsize]
@Test
void virementAccepte_lesSoldesSontMisAJour() {
  page.navigate(url() + "/");

  page.getByTestId("virement-source").selectOption("FR76-001");
  page.getByTestId("virement-cible").selectOption("FR76-002");
  page.getByLabel("Montant (€)").fill("150");
  page.getByRole(AriaRole.BUTTON,
      new Page.GetByRoleOptions().setName("Virer")).click();

  assertThat(page.getByTestId("message"))
      .hasText("Virement effectué");
  assertThat(page.getByTestId("solde-FR76-001"))
      .hasText("350.00 €");
}
\end{lstlisting}
    \end{column}
    \begin{column}{0.33\textwidth}
      \small
      \begin{itemize}
        \item \texttt{page} : un onglet neuf pour chaque test
        \item \texttt{getBy\ldots} : trouver un élément (un \textit{locator})
        \item \texttt{selectOption}, \texttt{fill}, \texttt{click} : agir comme l'utilisateur
        \item \texttt{assertThat(locator)} : assertions de Playwright, qui \textbf{attendent} que la condition soit vraie
        \item L'application tourne sur un port aléatoire (\texttt{RANDOM\_PORT})
      \end{itemize}
    \end{column}
  \end{columns}
\end{frame}

\begin{frame}[fragile]{Trouver les éléments : des sélecteurs robustes}
  \begin{columns}[T]
    \begin{column}{0.5\textwidth}
      \begin{alertblock}{Fragile : lié à la structure de la page}
\begin{lstlisting}[basicstyle=\lstcode\scriptsize]
page.locator("form select:nth-child(2)")
page.locator("div > div:nth-child(3) button")
page.locator("//table/tbody/tr[1]/td[3]")
page.locator(".btn.btn-primary.mt-2")
\end{lstlisting}
        \small On ajoute un champ, on change une classe CSS : tout casse.
      \end{alertblock}
    \end{column}
    \begin{column}{0.46\textwidth}
      \begin{exampleblock}{Robuste : lié à ce que voit l'utilisateur}
\begin{lstlisting}[basicstyle=\lstcode\scriptsize]
page.getByRole(AriaRole.BUTTON,
    new GetByRoleOptions().setName("Virer"))
page.getByLabel("Montant (€)")
page.getByText("Solde insuffisant")
page.getByTestId("solde-FR76-001")
\end{lstlisting}
      \end{exampleblock}
    \end{column}
  \end{columns}
  \vspace{2mm}
  \small
  \begin{itemize}
    \item Ordre de préférence : \textbf{rôle} et nom accessible, puis \textbf{libellé}, \textbf{texte}, et en dernier recours un attribut \texttt{data-testid} ajouté exprès pour les tests
    \item Bonus : si \texttt{getByRole} ne trouve pas le bouton, c'est souvent que la page est \textbf{mal accessible} (un \texttt{<div>} cliquable au lieu d'un \texttt{<button>})
  \end{itemize}
\end{frame}

\begin{frame}[fragile]{Attendre correctement}
  \begin{columns}[T]
    \begin{column}{0.5\textwidth}
      \begin{alertblock}{À éviter}
\begin{lstlisting}[basicstyle=\lstcode\scriptsize]
bouton.click();
// lu tout de suite : le serveur n'a pas répondu
String msg = page.locator("#message").textContent();
assertThat(msg).isEqualTo("Virement effectué");

bouton.click();
Thread.sleep(2000);   // assez ? trop ? sur la CI ?
\end{lstlisting}
      \end{alertblock}
    \end{column}
    \begin{column}{0.46\textwidth}
      \begin{exampleblock}{À faire}
\begin{lstlisting}[basicstyle=\lstcode\scriptsize]
bouton.click();
assertThat(page.getByTestId("message"))
    .hasText("Virement effectué");
\end{lstlisting}
      \end{exampleblock}
      \small
      \begin{itemize}
        \item Les \textbf{actions} attendent que l'élément soit visible, activé, stable
        \item Les \textbf{assertions} \texttt{assertThat(locator)} réessaient jusqu'à ce que la condition soit vraie (5 s par défaut)
        \item \texttt{Thread.sleep} : trop court, le test est instable ; trop long, la suite est lente
      \end{itemize}
    \end{column}
  \end{columns}
\end{frame}

\begin{frame}[fragile]{Le patron Page Object}
  \begin{columns}[T]
    \begin{column}{0.62\textwidth}
\begin{lstlisting}[basicstyle=\lstcode\scriptsize]
public class PageComptes {
  private final Page page;

  public PageComptes virer(String source, String cible, String montant) {
    page.getByTestId("virement-source").selectOption(source);
    page.getByTestId("virement-cible").selectOption(cible);
    page.getByLabel("Montant (€)").fill(montant);
    page.getByRole(BUTTON, new GetByRoleOptions().setName("Virer")).click();
    return this;
  }
  public Locator message() { return page.getByTestId("message"); }
  public Locator solde(String numero) { ... }
}

// dans le test
comptes.virer("FR76-001", "FR76-002", "150");
assertThat(comptes.message()).hasText("Virement effectué");
\end{lstlisting}
    \end{column}
    \begin{column}{0.35\textwidth}
      \small
      \begin{itemize}
        \item Une classe par écran (ou composant), qui \textbf{cache} les sélecteurs et expose des actions métier
        \item Les tests se lisent comme un scénario
        \item La page change ? On corrige \textbf{une} classe, pas trente tests
        \item Pas d'assertions dans le Page Object : il décrit l'écran, le test décide de ce qui est attendu
      \end{itemize}
    \end{column}
  \end{columns}
\end{frame}

\begin{frame}{Données et isolation}
  \begin{itemize}
    \item Les mêmes règles qu'au CM4, en plus difficile : le navigateur et le serveur sont deux processus, \textbf{pas de rollback} automatique
    \item Avant chaque test, remettre l'application dans un état connu :
      \begin{itemize}
        \item base de test neuve (Testcontainers) et jeu de données chargé avant chaque test
        \item ou un point d'accès réservé aux tests (\texttt{POST /api/test/reinitialiser}), désactivé en production
      \end{itemize}
    \item Chaque test crée ou utilise \textbf{ses} comptes : deux tests qui virent depuis le même compte se gênent si on les exécute en parallèle
    \item Un nouvel onglet propre (\texttt{BrowserContext}) par test : pas de cookies ni de session hérités
    \item Ne pas passer par l'interface pour préparer les données (lent) : utiliser l'API ou la base, et réserver l'interface à l'action testée
  \end{itemize}
\end{frame}

\begin{frame}{Déboguer un test d'interface}
  \begin{columns}[T]
    \begin{column}{0.48\textwidth}
      \begin{itemize}
        \item \textbf{Mode visible} (\textit{headed}) et ralenti (\texttt{setSlowMo}) : on regarde le test jouer
        \item \textbf{Générateur de code} (\texttt{codegen}) : on clique dans la page, Playwright écrit le code Java
          \begin{itemize}
            \item bon point de départ\ldots{} à nettoyer (Page Object, assertions)
          \end{itemize}
        \item Captures d'écran et vidéos en cas d'échec
      \end{itemize}
    \end{column}
    \begin{column}{0.48\textwidth}
      \begin{block}{Le \textit{trace viewer}}
        \begin{itemize}
          \item Une trace (\texttt{.zip}) enregistre chaque action, une capture du DOM avant et après, la console et les requêtes réseau
          \item On la rejoue pas à pas, après coup
          \item Indispensable quand le test échoue \textbf{en CI}, sur une machine qu'on ne voit pas (CM6)
        \end{itemize}
      \end{block}
    \end{column}
  \end{columns}
  \vspace{3mm}
  \begin{center}
    \small En CI, le navigateur tourne sans fenêtre (\textit{headless}) : c'est le mode par défaut.
  \end{center}
\end{frame}

\begin{frame}[fragile]{Retour au CM2 : le même scénario, dans le navigateur}
  \begin{columns}[T]
    \begin{column}{0.44\textwidth}
\begin{lstlisting}[language={}, basicstyle=\lstcode\scriptsize]
Scénario: virement accepté
  Étant donné un compte "FR76-001"
    avec un solde de 500,00 €
  Quand le client vire 150,00 €
    de "FR76-001" vers "FR76-002"
  Alors le solde de "FR76-001"
    est de 350,00 €
\end{lstlisting}
    \end{column}
    \begin{column}{0.53\textwidth}
\begin{lstlisting}[basicstyle=\lstcode\scriptsize]
@When("le client vire {bigdecimal} € de {string} vers {string}")
public void vire(BigDecimal montant, String source, String cible) {
    pageComptes.virer(source, cible, montant.toPlainString());
}

@Then("le solde de {string} est de {bigdecimal} €")
public void soldeEst(String numero, BigDecimal attendu) {
    assertThat(pageComptes.solde(numero))
        .hasText(attendu.setScale(2) + " €");
}
\end{lstlisting}
    \end{column}
  \end{columns}
  \vspace{2mm}
  \small
  \begin{itemize}
    \item Le fichier \texttt{.feature} ne change pas : seules les \textit{step definitions} passent par l'interface, via le Page Object
    \item On choisit les scénarios qui méritent ce niveau : en général, un ou deux par fonctionnalité
  \end{itemize}
\end{frame}

\rdemo{1}{Un test qui pilote le navigateur}{%
  \item Le test joue sous nos yeux (mode visible, ralenti)
  \item La version « fragile » : sélecteurs \texttt{nth-child}, texte lu trop tôt\ldots{} rouge, alors que l'application marche
  \item Le \textit{trace viewer} : rejouer le test pas à pas, voir l'écran au moment de l'échec
  \item \textit{Si le temps le permet} : \texttt{codegen} écrit le test pendant qu'on clique
}{\texttt{VirementE2ETest}\\\texttt{FragileE2ETest}\\\texttt{pages/PageComptes}\\[1mm]\texttt{target/traces/*.zip}}

% ================================================================
\section{Ouverture : utilisabilité et accessibilité}

\begin{frame}{Tester l'utilisabilité}
  \begin{columns}[T]
    \begin{column}{0.5\textwidth}
      \begin{itemize}
        \item Un test automatique vérifie que le virement \textbf{fonctionne} ; il ne dit pas si un client \textbf{réussit} à le faire
        \item \textbf{Test utilisateur} : on observe de vrais utilisateurs réaliser des tâches (« virez 50 \euro{} à votre fils »)
          \begin{itemize}
            \item on mesure : réussite, temps, erreurs, hésitations
            \item ils pensent à voix haute ; on n'aide pas
            \item cinq utilisateurs révèlent déjà la plupart des gros problèmes (Nielsen)
          \end{itemize}
      \end{itemize}
    \end{column}
    \begin{column}{0.46\textwidth}
      \begin{block}{Quelques heuristiques de Nielsen}
        \small
        \begin{itemize}
          \item Rendre visible l'état du système (« virement en cours\ldots »)
          \item Parler la langue de l'utilisateur (pas « erreur 422 »)
          \item Prévenir les erreurs (confirmation au-delà de 1\,000 \euro{})
          \item Permettre d'annuler
          \item Aider à comprendre et corriger une erreur
        \end{itemize}
      \end{block}
    \end{column}
  \end{columns}
\end{frame}

\begin{frame}{L'accessibilité}
  \begin{itemize}
    \item Utilisable par tous : lecteur d'écran, clavier seul, malvoyance, daltonisme, troubles cognitifs\ldots
    \item Une \textbf{obligation légale} :
      \begin{itemize}
        \item référentiel français \textbf{RGAA} (fondé sur les WCAG internationales)
        \item depuis juin 2025, la directive européenne sur l'accessibilité s'applique aussi à des services privés, dont les \textbf{services bancaires} en ligne
      \end{itemize}
    \item Ce qui s'automatise : contrastes, textes alternatifs, libellés de champs, structure des titres
      \begin{itemize}
        \item outils : axe-core (utilisable depuis Playwright), Lighthouse dans le navigateur
        \item ils ne détectent qu'une \textbf{partie} des problèmes : le reste demande des tests manuels (clavier, lecteur d'écran)
      \end{itemize}
    \item Les sélecteurs \texttt{getByRole} et \texttt{getByLabel} ne fonctionnent que sur une page correctement balisée : de bons tests poussent vers une meilleure accessibilité
  \end{itemize}
\end{frame}

\begin{frame}{À retenir}
  \small
  \begin{itemize}
    \item Un test de bout en bout pilote un \textbf{vrai navigateur} sur l'application entière : réaliste, mais lent, fragile et coûteux à diagnostiquer
    \item Peu nombreux : les \textbf{parcours critiques} ; les règles métier restent testées plus bas
    \item Un test instable est à réparer en priorité : sinon plus personne ne croit aux tests
    \item Playwright : \texttt{getByRole}, \texttt{getByLabel}, \texttt{getByTestId} plutôt que des sélecteurs CSS ou XPath liés à la structure
    \item Jamais de \texttt{Thread.sleep} : actions et assertions \texttt{assertThat(locator)} \textbf{attendent} toutes seules
    \item \textbf{Page Object} : une classe par écran, qui cache les sélecteurs
    \item Données remises dans un état connu avant chaque test ; un contexte de navigateur neuf par test
    \item Mode visible, \texttt{codegen} et \textit{trace viewer} pour écrire et déboguer
    \item Fonctionner ne suffit pas : utilisabilité et accessibilité se testent aussi
  \end{itemize}
  \vspace{1mm}
  \begin{exampleblock}{TP5}
    Tester le front Vue du projet fil rouge avec Playwright.
  \end{exampleblock}
\end{frame}

\rquestionslide

\end{document}
